\documentclass[12pt,a4paper]{report}
\usepackage[utf8]{inputenc}
\usepackage[T1]{fontenc}
\usepackage[french]{babel}
\usepackage[margin=2.5cm]{geometry}
\usepackage{graphicx}
\usepackage{fancyhdr}
\usepackage{tikz}
\usepackage{listings}
\usepackage{xcolor}
\usepackage{hyperref}
\usepackage{array}
\usepackage{booktabs}

\lstset{
    language=Java,
    basicstyle=\ttfamily\small,
    keywordstyle=\color{blue}\textbf,
    commentstyle=\color{gray},
    stringstyle=\color{red},
    breaklines=true,
    showstringspaces=false,
    tabsize=4,
    backgroundcolor=\color{lightgray!20}
}

\title{\textbf{Rapport du Projet Web Service SOAP} \\ \textit{Gestion des Étudiants}}
\author{\textbf{Développé par:} Abdel \\ \textit{Université/École}}
\date{\today}

\pagestyle{fancy}
\fancyhf{}
\rhead{Web Service SOAP - Rapport}
\lhead{Projet Étudiant}
\cfoot{\thepage}

\begin{document}

% ===== PAGE DE TITRE =====
\begin{titlepage}
\centering
\vspace*{\fill}

{\LARGE\textbf{RAPPORT TECHNIQUE}}\\[0.5cm]
{\Large Web Service SOAP - Gestion des Étudiants}\\[1cm]

\rule{\linewidth}{2pt}\\[1cm]

{\large Développé par : \textbf{Abdel}}\\[0.5cm]
{\large Date : \textbf{\today}}\\[0.5cm]
{\large Technologie : \textbf{Spring Boot 3.2.4 + MySQL + SOAP}}\\[2cm]

\vspace*{\fill}

{\large \textit{Ce rapport présente l'architecture, la conception et l'implémentation} \\ 
\textit{d'un Web Service SOAP pour la gestion des données d'étudiants.}}

\end{titlepage}

\newpage
% ===== TABLE DES MATIÈRES =====
\tableofcontents
\newpage

% ===== CHAPITRE 1 =====
\chapter{Introduction}

\section{Contexte du Projet}
Ce projet met en place un Web Service SOAP (Simple Object Access Protocol) pour la gestion d'étudiants. L'application utilise les technologies modernes de développement Java avec le framework Spring Boot, couplée à une base de données MySQL pour la persistance des données.

\section{Objectifs}
\begin{itemize}
    \item Créer un service Web SOAP fonctionnel
    \item Implémenter les opérations CRUD (Create, Read, Update, Delete)
    \item Assurer la persistance des données en base MySQL
    \item Fournir une architecture scalable et maintenable
\end{itemize}

\section{Spécifications Techniques}
\begin{table}[h]
\centering
\begin{tabular}{|l|c|}
\hline
\textbf{Composant} & \textbf{Version/Valeur} \\
\hline
Java & 21 (LTS) \\
Spring Boot & 3.2.4 \\
Base de Données & MySQL 8.0+ \\
Port d'écoute & 8081 \\
Protocole & SOAP 1.1 \\
Namespace & http://kholty.com/student \\
\hline
\end{tabular}
\caption{Spécifications Techniques}
\end{table}

\newpage

% ===== CHAPITRE 2 =====
\chapter{Architecture et Conception}

\section{Architecture Générale}

\begin{figure}[h]
\centering
\begin{tikzpicture}[scale=1.0]
    % Client
    \draw[fill=blue!20] (0,8) rectangle (2,9) node[midway] {\textbf{Client}};
    \draw[->,thick] (1,8) -- (1,7);
    
    % SOAP Request
    \node[text=red,font=\small] at (2.5,7.5) {SOAP Request};
    
    % Web Service Layer
    \draw[fill=green!20] (0,5.5) rectangle (2,7) node[midway] {\textbf{SOAP Endpoint} \\ \texttt{/ws}};
    \draw[->,thick] (1,5.5) -- (1,4.5);
    
    % Business Logic
    \draw[fill=yellow!20] (0,3) rectangle (2,4.5) node[midway] {\textbf{Endpoint} \\ \texttt{StudentEndpoint}};
    \draw[->,thick] (1,3) -- (1,2);
    
    % Repository Layer
    \draw[fill=orange!20] (0,0.5) rectangle (2,2) node[midway] {\textbf{Repository} \\ \texttt{StudentRepository}};
    \draw[->,thick] (2.5,1.25) -- (4,1.25);
    
    % Database
    \draw[fill=red!20] (4,0.5) rectangle (6,2) node[midway] {\textbf{MySQL Database} \\ \texttt{soap.student}};
    
    % Response Path
    \draw[dashed,->,thick,blue] (6,1.25) -- (5,1.25);
    \draw[dashed,->,thick,blue] (2.5,1.25) -- (2,1.25);
    \draw[dashed,->,thick,blue] (1,2) -- (1,3);
    \draw[dashed,->,thick,blue] (1,4.5) -- (1,5.5);
    \node[text=blue,font=\small] at (2.5,2.5) {SOAP Response};
    \draw[dashed,->,thick,blue] (1,7) -- (1,8);
    
\end{tikzpicture}
\caption{Architecture du Web Service SOAP}
\end{figure}

\section{Flux de Données (Request/Response)}

\subsection{Cycle Complet d'une Requête}

\textbf{1. Requête CLIENT SOAP:}
\begin{lstlisting}
POST /ws HTTP/1.1
Host: localhost:8081
Content-Type: application/soap+xml

<Envelope xmlns="http://schemas.xmlsoap.org/soap/envelope/">
  <Body>
    <addStudentRequest xmlns="http://kholty.com/student">
      <name>Ali Ahmed</name>
      <email>ali@kholty.com</email>
    </addStudentRequest>
  </Body>
</Envelope>
\end{lstlisting}

\textbf{2. Traitement SERVER:}
\begin{enumerate}
    \item MessageDispatcherServlet reçoit la requête
    \item PayloadRootQNameEndpointMapping mappe à StudentEndpoint
    \item Méthode \texttt{addStudent()} est appelée
    \item Données extraites et validées
    \item StudentRepository.save() persiste en MySQL
    \item BusinessLogic retourne StudentType
\end{enumerate}

\textbf{3. Réponse SERVER SOAP:}
\begin{lstlisting}
HTTP/1.1 200 OK
Content-Type: application/soap+xml

<Envelope xmlns="http://schemas.xmlsoap.org/soap/envelope/">
  <Body>
    <addStudentResponse xmlns="http://kholty.com/student">
      <student>
        <id>1</id>
        <name>Ali Ahmed</name>
        <email>ali@kholty.com</email>
      </student>
    </addStudentResponse>
  </Body>
</Envelope>
\end{lstlisting}

\newpage

% ===== CHAPITRE 3 =====
\chapter{Implémentation Détaillée}

\section{Structure du Projet}

\begin{verbatim}
Web_service_SOAP/
├── src/main/java/
│   └── com/example/demo/
│       ├── WebServiceApplication.java (Point d'entrée)
│       ├── model/
│       │   └── Student.java (@Entity JPA)
│       ├── repository/
│       │   └── StudentRepository.java (JpaRepository)
│       ├── Exception/
        │   └── StudentException.java (@SoapFault)
│       └── web/services/soap/
│           ├── endpoint/
│           │   └── StudentEndpoint.java (SOAP Operations)
│           └── config/
│               └── WebServiceConfig.java (Configuration)
├── src/main/resources/
│   ├── student.xsd (Schéma XML)
│   └── application.properties (Configuration)
└── pom.xml (Dépendances)
\end{verbatim}

\section{Modèle de Données}

\subsection{Entité JPA - Student}
\begin{lstlisting}
@Entity
@Data
@NoArgsConstructor
@AllArgsConstructor
public class Student {
    @Id
    @GeneratedValue(strategy = GenerationType.IDENTITY)
    private Long id;
    private String name;
    private String email;
}
\end{lstlisting}

\subsection{Schéma XML (student.xsd)}
Le schéma XSD définit les contrats pour :
\begin{itemize}
    \item \textbf{addStudentRequest/Response} ✓ Implémenté
    \item \textbf{getStudentRequest/Response} ✓ Implémenté
    \item \textbf{getAllStudentsRequest/Response} (Défini en XSD, À implémenter)
\end{itemize}

\section{Gestion des Exceptions}

\subsection{StudentException - Exception SOAP personnalisée}
\begin{lstlisting}
@SoapFault(faultCode = FaultCode.CUSTOM,
    customFaultCode = "{http://kholty.com/Student}001_not found")
public class StudentException extends RuntimeException {
    public StudentException(String message) {
        super(message);
    }
}
\end{lstlisting}

L'annotation \texttt{@SoapFault} transforme les exceptions Java en réponses SOAP Fault standard, permettant au client de recevoir une erreur SOAP bien formée au lieu d'une exception HTTP 500.

\section{Operations SOAP}

\subsection{1. Add Student (CREATE)}
\begin{lstlisting}
@Transactional
@PayloadRoot(namespace = namespace, localPart = "addStudentRequest")
@ResponsePayload
public AddStudentResponse addStudent(@RequestPayload AddStudentRequest request) {
    Student student = new Student();
    student.setName(request.getName());
    student.setEmail(request.getEmail());
    
    Studentrepo.save(student);  // Persiste en MySQL
    
    // Retour
    StudentType st = new StudentType();
    st.setId(student.getId());
    st.setName(student.getName());
    st.setEmail(student.getEmail());
    
    AddStudentResponse response = new AddStudentResponse();
    response.setStudent(st);
    return response;
}
\end{lstlisting}

\subsection{2. Get Student (READ)}
\begin{lstlisting}
@PayloadRoot(namespace = namespace, localPart = "getStudentRequest")
@ResponsePayload
public GetStudentResponse getStudentById(@RequestPayload GetStudentRequest request) 
    throws StudentException {
    
    Optional<Student> student = Studentrepo.findById(request.getId());
    Student st = new Student();
    
    if (student.isPresent()) {
        st = student.get();
    } else {
        throw new StudentException("student with id ( " + 
            request.getId() + " ) not found!");
    }
    
    // Mapper vers StudentType
    StudentType stf = new StudentType();
    stf.setId(st.getId());
    stf.setName(st.getName());
    stf.setEmail(st.getEmail());
    
    GetStudentResponse response = new GetStudentResponse();
    response.setStudent(stf);
    return response;
}
\end{lstlisting}

\newpagapplication.name=demo
server.port=8081

# HikariCP settings (Default connection pool)
spring.datasource.hikari.minimumIdle=5
spring.datasource.hikari.maximumPoolSize=20
spring.datasource.hikari.idleTimeout=30000
spring.datasource.hikari.maxLifetime=2000000
spring.datasource.hikari.connectionTimeout=30000
spring.datasource.hikari.poolName=HikariPoolBooks

spring.jpa.hibernate.ddl-auto=update
spring.jpa.show-sql=true

spring.datasource.url=jdbc:mysql://${MYSQL_HOST:localhost}:3306/soap
spring.datasource.username=root
spring.datasource.password=kholtyySQL}

\begin{lstlisting}[language=properties]
# application.properties
spring.datasource.url=jdbc:mysql://localhost:3306/soap
spring.datasource.username=root
spring.datasource.password=kholty

spring.jpa.hibernate.ddl-auto=update
spring.jpa.show-sql=true
\end{lstlisting}

\section{Diagramme ER - Base de Données}

\begin{figure}[h]
\centering
\begin{tikzpicture}
    \draw[fill=blue!20] (2,2) rectangle (5,3.5);
    \node[text centered] at (3.5,3.2) {\textbf{STUDENT}};
    \node[font=\small] at (2.2,2.8) {id (PK)};
    \node[font=\small] at (2.2,2.4) {name (VARCHAR)};
    \node[font=\small] at (2.2,2.0) {email (VARCHAR)};
\end{tikzpicture}
\caption{Modèle de Données MySQL}
\end{figure}

\section{Opérations CRUD en Base}

\begin{table}[h]
\centering
\begin{tabular}{|c|c|c|}
\hline
\textbf{Opération} & \textbf{Méthode} & \textbf{SQL} \\
\hline
CREATE & \texttt{save()} & \texttt{INSERT INTO student ...} \\
READ & \texttt{findById()} & \texttt{SELECT * FROM student WHERE id=?} \\
UPDATE & \texttt{save()} & \texttt{UPDATE student SET ...} \\
DELETE & \texttt{deleteById()} & \texttt{DELETE FROM student WHERE id=?} \\
\hline
\end{tabular}
\caption{Opérations CRUD}
\end{table}

\section{Exemple de Cycle Complet}

\begin{enumerate}
    \item \textbf{SOAP Request reçue:} addStudentRequest avec name="Ali", email="ali@test.com"
    \item \textbf{StudentEndpoint traite:} Crée un nouvel objet Student
    \item \textbf{StudentRepository.save():} Exécute \texttt{INSERT INTO student (name, email) VALUES ('Ali', 'ali@test.com')}
Ce Web Service SOAP fournit une solution fonctionnelle pour la gestion des étudiants avec :
\begin{itemize}
    \item ✓ Architecture en couches bien définies (Endpoint → Repository → JPA/MySQL)
    \item ✓ Persistance MySQL fiable via HikariCP (pool 5-20 connexions)
    \item ✓ Transactions SOAP atomiques avec \texttt{@Transactional}
    \item ✓ Gestion d'erreurs avec \texttt{@SoapFault} personnalisée
    \item ✓ Configuration sur port 8081 (évite les conflits)
    \item ✓ WSDL généré automatiquement via DefaultWsdl11Definition
    \item ✓ Deux opérations SOAP complètes : addStudent, getStudentById
% ===== CHAPITRE 5 =====
\chapter{Conclusion et Améliorations Futures}

\section{Résumé}
Ce Web Service SOAP fournit une solution robuste pour la gestion des étudiants avec :
\begin{itemize}
    \item Implémenter l'opération getAllStudents (définie en XSD mais pas encore codée)
    \item Implémenter l'opération deleteStudent
    \item Ajouter Spring Security pour authentification/autorisation
    \item Ajouter des validations (format email, longueur min/max, etc.)
    \item Implémenter des logs applicatifs structurés
    \item Ajouter des tests unitaires JUnit 5 et intégration
    \item ✓ Port sécurisé 8081
\end{itemize}

\section{Améliorations Recommandées}

\subsection{Court Terme}
\begin{itemize}
    \item Ajouter Spring Security pour authentification
    \item Implémenter les opérations getAllStudents et deleteStudent
    \item Ajouter des validations (email format, etc.)
    \item Implémenter des logs applicatifs
\end{itemize}

\subsection{Long Terme}
\begin{itemize}
    \item Migrer vers REST API (plus moderne que SOAP)
    \item Ajouter caching (Redis)
    \item Implémenter CI/CD (Jenkins, GitHub Actions)
    \item Conteneuriset Outils}

\subsection{Dépendances Maven}
\begin{itemize}
    \item \textbf{spring-boot-starter-web-services:} Support SOAP/WS et MessageDispatcherServlet
    \item \textbf{spring-boot-starter-data-jpa:} Persistence JPA/Hibernate
    \item \textbf{mysql-connector-java:} Driver MySQL pour connexion base de données
    \item \textbf{lombok:} Annotations @Data, @NoArgsConstructor, @AllArgsConstructor
    \item \textbf{jakarta.transaction:} Support @Transactional (Jakarta EE)
\end{itemize}

\subsection{Plugins Maven}
\begin{itemize}
    \item \textbf{jaxb2-maven-plugin:} Génération classes SOAP depuis XSD (addStudentRequest, StudentType, etc.)
    \item \textbf{spring-boot-maven-plugin:} Build JAR exécutable et hot reload
    \item \textbf{maven-compiler-plugin:} Compilation Java 21
\end{itemize}

\subsection{Infrastructure}
\begin{itemize}
    \item \textbf{HikariCP:} Pool de connexions MySQL (5-20 connexions)
    \item \textbf{MySQL 8.0+:} Base de données avec schéma "soap" et table "student"
    \item \textbf{Spring Boot Embedded Tomcat:} Serveur HTTP sur port 8081/WS
    \item \textbf{spring-boot-starter-data-jpa:} Persistence
    \item \textbf{mysql-connector-java:} Driver MySQL
    \item \textbf{lombok:} Annotations @Data, @NoArgsConstructor
    \item \textbf{jaxb2-maven-plugin:} Génération code depuis XSD
\end{itemize}

\vspace{2cm}

\section{État Actuel du Projet}

\subsection{✓ Fonctionnalités Implémentées et Testées}
\begin{enumerate}
    \item \textbf{Add Student (CREATE):} Crée un étudiant en base MySQL, génère ID auto-incrémenté
    \item \textbf{Get Student (READ):} Récupère un étudiant par ID avec gestion d'erreur StudentException
    \item \textbf{Configuration SOAP:} MessageDispatcherServlet, PayloadRootQNameEndpointMapping, génération WSDL
    \item \textbf{Persistance MySQL:} HikariCP pool, JPA/Hibernate, auto-migration du schéma
    \item \textbf{Port 8081:} Application démarre sans conflit, logs Tomcat vérifiés
    \item \textbf{Gestion erreurs:} @SoapFault personnalisée pour erreurs métier (NOT_FOUND)
\end{enumerate}

\subsection{⏳ À Implémenter}
\begin{enumerate}
    \item \textbf{Get All Students (READ LIST):} getAllStudentsRequest/Response définis en XSD, pas encore d'endpoint
    \item \textbf{Delete Student:} Opération non présente en XSD ni en code
    \item \textbf{Update Student:} Opération non présente en XSD ni en code
    \item \textbf{Spring Security:} Authentification/autorisation non yet configurée
    \item \textbf{Validation:} Annotations @Valid, @NotNull, @Email manquantes
    \item \textbf{Logging:} Pas de logs structurés (SLF4J/Logback)
\end{enumerate}

\subsection{Notes Techniques}
\begin{itemize}
    \item \textbf{Namespace SOAP:} http://kholty.com/student (doit correspondre dans @PayloadRoot)
    \item \textbf{Annotations critiques:} \texttt{@Endpoint}, \texttt{@PayloadRoot}, \texttt{@ResponsePayload}
    \item \textbf{Cas sensible:} localPart="addStudentRequest" (minuscule) doit matcher l'élément XSD
    \item \textbf{Optional API:} Toujours checker \texttt{isPresent()} avant \texttt{get()}
    \item \textbf{Transactions:} \texttt{@Transactional} sur addStudent pour atomicité
\end{itemize}

\centerline{\textit{Rapport certifié exact pour version git commit : 2026-03-12-rapport-v2}}

\end{document}
